\section{Cơ sở lý thuyết}

\subsection{Chuyển động ném xiên -- Mô hình lý tưởng}
\subsubsection{Phát biểu bài toán}
Xét một chất điểm khối lượng $m$ được ném từ gốc tọa độ $O$ với vận tốc ban đầu $v_0$ hợp với phương ngang một góc $\alpha$. Trong trường hợp lý tưởng (bỏ qua lực cản không khí), chất điểm chỉ chịu tác dụng của trọng lực $\vec{P} = m\vec{g}$, và quỹ đạo là một parabol đối xứng.

Tuy nhiên, trong thực tế, khi vật chuyển động trong môi trường chất lưu (không khí, nước,...), vật luôn chịu tác dụng của \textbf{lực cản môi trường}. Lực cản này phụ thuộc vào vận tốc và làm thay đổi đáng kể quỹ đạo so với mô hình lý tưởng.

\subsubsection{Phương trình chuyển động lý tưởng}
Chọn hệ tọa độ $Oxy$ với $Ox$ theo phương ngang, $Oy$ theo phương thẳng đứng hướng lên. Khi bỏ qua lực cản, phương trình Newton II cho chuyển động ném xiên:
$$m\vec{a} = m\vec{g}$$

Chiếu lên hai phương:
\begin{equation}
\begin{cases}
x(t) = v_0 \cos\alpha \cdot t \\
y(t) = v_0 \sin\alpha \cdot t - \dfrac{1}{2}gt^2
\end{cases}
\label{eq:ideal}
\end{equation}

Trong trường hợp này, quỹ đạo là parabol đối xứng, tầm xa cực đại đạt được khi $\alpha = 45°$:
$$R_{max} = \frac{v_0^2}{g}$$


\subsection{Mô hình có lực cản môi trường}
\subsubsection{Lực cản tuyến tính theo vận tốc}
Khi vật chuyển động trong môi trường nhớt với vận tốc không quá lớn, lực cản có dạng tuyến tính theo vận tốc:
$$\vec{F}_{cản} = -h\vec{v}$$

trong đó $h$ (kg/s) là hệ số cản của môi trường. Hệ số này phụ thuộc vào hình dạng vật, độ nhớt của môi trường và tiết diện ngang của vật thể.

Phương trình Newton II tổng quát cho bài toán trở thành:
\begin{equation}
m\vec{a} = m\vec{g} - h\vec{v}
\label{eq:newton}
\end{equation}

\subsubsection{Hệ phương trình vi phân theo hai phương}
Chiếu phương trình (\ref{eq:newton}) lên hệ tọa độ $Oxy$:

\textbf{Phương Ox (ngang):}
\begin{equation}
m\ddot{x} = -h\dot{x} \quad \Leftrightarrow \quad m\ddot{x} + h\dot{x} = 0
\label{eq:ode_x}
\end{equation}

\textbf{Phương Oy (đứng):}
\begin{equation}
m\ddot{y} = -mg - h\dot{y} \quad \Leftrightarrow \quad m\ddot{y} + h\dot{y} = -mg
\label{eq:ode_y}
\end{equation}

Đây là hệ hai phương trình vi phân cấp hai, tuyến tính, hệ số hằng. Phương trình (\ref{eq:ode_x}) là ODE thuần nhất, phương trình (\ref{eq:ode_y}) là ODE không thuần nhất.

\subsubsection{Điều kiện ban đầu}
\begin{equation}
\begin{cases}
x(0) = 0, \quad y(0) = 0 \\
\dot{x}(0) = v_{0x} = v_0 \cos\alpha \\
\dot{y}(0) = v_{0y} = v_0 \sin\alpha
\end{cases}
\label{eq:ic}
\end{equation}

\subsection{Nghiệm giải tích}
Đặt $k = h/m$ là hệ số cản riêng (đơn vị: $s^{-1}$).

\subsubsection{Nghiệm theo phương Ox}
Phương trình (\ref{eq:ode_x}) là ODE tuyến tính thuần nhất bậc 2 với hệ số hằng. Phương trình đặc trưng $r^2 + kr = 0$ cho nghiệm $r_1 = 0$, $r_2 = -k$. Áp dụng điều kiện ban đầu:

\textbf{Vận tốc:}
\begin{equation}
v_x(t) = v_0 \cos\alpha \cdot e^{-kt}
\label{eq:vx}
\end{equation}

\textbf{Vị trí:}
\begin{equation}
x(t) = \frac{m \cdot v_0 \cos\alpha}{h}\left(1 - e^{-\frac{h}{m}t}\right)
\label{eq:x}
\end{equation}

Nhận xét: Khi $t \to \infty$, $v_x \to 0$ và $x \to \dfrac{mv_0\cos\alpha}{h}$ (giá trị hữu hạn). Đây là điểm khác biệt cơ bản so với trường hợp không cản (vận tốc ngang không đổi).

\subsubsection{Nghiệm theo phương Oy}
Phương trình (\ref{eq:ode_y}) là ODE tuyến tính không thuần nhất. Nghiệm gồm nghiệm thuần nhất và nghiệm riêng (nghiệm dừng $v_y = -g/k$):

\textbf{Vận tốc:}
\begin{equation}
v_y(t) = \left(v_0\sin\alpha + \frac{mg}{h}\right)e^{-\frac{h}{m}t} - \frac{mg}{h}
\label{eq:vy}
\end{equation}

\textbf{Vị trí:}
\begin{equation}
y(t) = \frac{m}{h^2}\left[\left(g \cdot m + h \cdot v_0\sin\alpha\right)\left(1 - e^{-\frac{h}{m}t}\right) - g \cdot h \cdot t\right]
\label{eq:y}
\end{equation}

\subsection{Các đại lượng đặc trưng}
\subsubsection{Vận tốc giới hạn (Terminal velocity)}
Khi $t \to \infty$, từ phương trình (\ref{eq:vy}):
$$v_\infty = -\frac{mg}{h}$$

Đây là vận tốc giới hạn khi lực cản cân bằng hoàn toàn với trọng lực. Vật rơi đều với tốc độ không đổi $|v_\infty| = mg/h$.

\subsubsection{Thời gian đạt độ cao cực đại}
Tại đỉnh quỹ đạo, $v_y = 0$. Từ (\ref{eq:vy}):
\begin{equation}
t_{max} = \frac{m}{h}\ln\left(1 + \frac{h v_0\sin\alpha}{mg}\right)
\label{eq:tmax}
\end{equation}

So với trường hợp không cản ($t_{max,0} = v_0\sin\alpha / g$), thời gian leo lên \textbf{ngắn hơn} do lực cản hướng ngược vận tốc làm giảm tốc mạnh hơn.

\subsubsection{Năng lượng}
Cơ năng của hệ ($E = KE + PE$) không bảo toàn do lực cản là lực không bảo toàn:
\begin{itemize}
    \item Động năng: $KE = \frac{1}{2}m(v_x^2 + v_y^2)$
    \item Thế năng: $PE = mgy$
    \item Năng lượng mất do cản: $W_{cản} = E_0 - (KE + PE)$
\end{itemize}

Định luật bảo toàn năng lượng tổng quát: $E_0 = KE + PE + W_{cản}$ luôn được thỏa mãn tại mọi thời điểm.
